\title{Chain-of-Thought Prompting Elicits Reasoning in Large Language Models}

\begin{abstract}
We investigate how producing a \textit{chain of thought}—a sequence of intermediate reasoning steps—can substantially enhance the capacity of large language models to handle complex reasoning tasks. Specifically, we demonstrate that such reasoning capabilities emerge naturally in sufficiently large language models through a straightforward technique called \textit{chain-of-thought prompting}, where a few chain-of-thought examples are included as prompts.

Tests conducted on three large language models reveal that chain-of-thought prompting boosts performance across various arithmetic, commonsense, and symbolic reasoning challenges. The observed improvements can be remarkable. For example, prompting a PaLM 540B model with just eight chain-of-thought exemplars attains state-of-the-art results on the GSM8K benchmark for math word problems, outperforming even finetuned GPT-3 combined with a verifier.
\end{abstract}

\section{Introduction}
Recent advances in NLP have been driven by language models \citep[][\textit{inter alia}]{peters-etal-2018-deep,devlin-etal-2019-bert,brown2020language}. Increasing the scale of language models has been demonstrated to yield numerous advantages, including enhanced accuracy and improved sample efficiency \citep[\textit{inter alia}]{kaplan2020scaling,brown2020language}. Nonetheless, simply enlarging the model size has not sufficed to achieve strong performance on demanding tasks such as arithmetic, commonsense, and symbolic reasoning \citep{rae2021scaling}.

This study examines how to unlock reasoning capabilities in large language models by means of a simple approach inspired by two insights. First, arithmetic reasoning methods can benefit from generating natural language explanations that lead to the answer. Previous research has enabled models to produce natural language intermediate steps either by training from scratch \citep{ling-etal-2017-program} or by finetuning a pretrained model \citep{cobbe2021training}, alongside neuro-symbolic approaches that employ formal languages instead of natural language \citep{roy-roth-2015-solving, chiang-chen-2019-semantically,amini-etal-2019-mathqa, chen2019neural}. Second, large language models present an exciting opportunity for in-context few-shot learning via \emph{prompting}. Rather than finetuning a separate model checkpoint for each new task, one can simply provide the model with a few input--output examples illustrating the task. Notably, this approach has proven effective for a variety of straightforward question-answering problems \citep{brown2020language}.

Nonetheless, both of the aforementioned concepts have significant drawbacks. For rationale-augmented training and finetuning techniques, generating a substantial collection of high-quality rationales is resource-intensive, as it is considerably more complex than the straightforward input--output pairs typically employed in conventional machine learning. Regarding the traditional few-shot prompting approach utilized in \citet{brown2020language}, it demonstrates poor performance on tasks necessitating reasoning capabilities and often fails to show significant improvement even as the language model scale increases \citep{rae2021scaling}. In this study, we integrate the advantages of these two strategies while circumventing their respective drawbacks. Specifically, we investigate the capacity of language models to execute few-shot prompting for reasoning tasks when given a prompt comprising triples: $\langle$input, \emph{chain of thought}, output$\rangle$.

A \emph{chain of thought} refers to a sequence of intermediate natural language reasoning steps that culminate in the final output, and we label this method as \textit{chain-of-thought prompting}. An example of such a prompt is displayed in \cref{fig:pull-figure}.

We conduct empirical assessments on benchmarks involving arithmetic, commonsense, and symbolic reasoning, demonstrating that chain-of-thought prompting surpasses standard prompting, occasionally by a remarkable margin. \cref{fig:pull-bar-chart} depicts one such outcome—on the GSM8K benchmark of math word problems \citep{cobbe2021training}, chain-of-thought prompting using \palm{} 540B significantly outperforms standard prompting and sets a new state-of-the-art. Employing a prompting-only strategy is crucial because it obviates the need for an extensive training dataset and allows a single model checkpoint to tackle multiple tasks without sacrificing generality. This research highlights how large language models can learn from a few examples using natural language task data (in contrast to automatically acquiring input-output patterns through a large training dataset).

\section{Chain-of-Thought Prompting} When tackling a complex reasoning problem, such as a multi-step math word problem, individuals often reflect on their own thought process. Typically, the problem is broken down into intermediate stages, solving each step before arriving at the final solution: \textit{``After Jane gives 2 flowers to her mom she has 10 $\ldots$ then after she gives 3 to her dad she will have 7 $\ldots$ so the answer is 7.''} This paper aims to equip language models with the capability to produce a similar \textit{chain of thought}—a logically connected sequence of intermediate reasoning steps that culminate in the final answer. We demonstrate that language models of sufficient size can generate such chains of thought when examples of chain-of-thought reasoning are included in few-shot prompting exemplars.

\cref{fig:pull-figure} illustrates a model generating a chain of thought to solve a math word problem, which it would have otherwise answered incorrectly. In this instance, the chain of thought resembles a solution and can be interpreted as one; however, we prefer to label it as a chain of thought to emphasize that it emulates a step-by-step reasoning process leading to the answer (also, solutions or explanations generally follow the final answer \citep[][\textit{inter alia}]{narang2020wt5,wiegreffe2021reframing,lampinen2022can}).

Chain-of-thought prompting offers several appealing advantages as a technique to enhance reasoning capabilities in language models.

\begin{enumerate}[topsep=1pt,itemsep=0ex]
    \item Firstly, chain of thought enables models to break down multi-step problems into smaller intermediate stages, allowing for more computational effort to be focused on problems that involve additional reasoning steps.
    \item Secondly, a chain of thought offers an interpretable insight into the model's behavior, indicating the possible reasoning process behind a specific answer and enabling the identification and correction of errors in the reasoning path (although fully understanding the computations a model uses to support an answer remains unresolved).
    \item Thirdly, chain-of-thought reasoning is applicable to tasks such as mathematical word problems, commonsense reasoning, and symbolic manipulation, and potentially to any problem that humans solve through language.
    \item Lastly, chain-of-thought reasoning can be easily induced in sufficiently large pre-trained language models by including examples of chain-of-thought sequences as part of the few-shot prompting exemplars.
\end{enumerate}

In our empirical evaluations, we will demonstrate the effectiveness of chain-of-thought prompting for arithmetic reasoning (\cref{sec:arithmetic-reasoning}), commonsense reasoning (\cref{sec:commonsense-reasoning}), and symbolic reasoning (\cref{sec:symbolic-reasoning}).

\section{Arithmetic Reasoning}\label{sec:arithmetic-reasoning}
We start by examining math word problems of the type shown in \cref{fig:pull-figure}, which assess the arithmetic reasoning capabilities of language models. Although straightforward for humans, arithmetic reasoning remains challenging for language models \citep[][\textit{inter alia}]{hendrycks2021measuring,patel-etal-2021-nlp}. Remarkably, chain-of-thought prompting combined with a 540B parameter language model achieves performance comparable to task-specific fine-tuned models across multiple benchmarks, even setting new state-of-the-art results on the difficult GSM8K benchmark \citep{cobbe2021training}.

\subsection{Experimental Setup}

We investigate chain-of-thought prompting across various language models and a range of benchmarks.

\textbf{Benchmarks.}
Our evaluation includes the following five math word problem datasets:
\textbf{(1)} the \textbf{GSM8K} benchmark \citep{cobbe2021training}, 
\textbf{(2)} the \textbf{SVAMP} dataset featuring math word problems with different structures \citep{patel-etal-2021-nlp},
\textbf{(3)} the \textbf{ASDiv} dataset containing diverse math word problems \citep{miao-etal-2020-diverse},
\textbf{(4)} the \textbf{AQuA} dataset of algebraic word problems, and
\textbf{(5)} the \textbf{MAWPS} benchmark \citep{koncel-kedziorski-etal-2016-mawps}. Sample problems are provided in Appendix \cref{tab:math-datasets}.

\textbf{Standard prompting.} As a baseline, we apply the standard few-shot prompting approach popularized by \citet{brown2020language}, where a language model receives in-context examples of input-output pairs prior to generating a response for a test example. The exemplars are structured as question-answer pairs. The model produces the final answer directly, as illustrated in \cref{fig:pull-figure} (left).

\textbf{Chain-of-thought prompting.} Our method involves enhancing each example in few-shot prompting by including a chain of thought that corresponds to the given answer, as depicted in \cref{fig:pull-figure} (right). Since most datasets only provide an evaluation split, we manually created a set of eight few-shot examples with chains of thought to use for prompting—one example with a chain of thought is shown in \cref{fig:pull-figure} (right), and the complete set is available in Appendix \cref{tab:appendix-mwp-prompts}. (These specific examples were not subjected to prompt engineering; robustness is analyzed in \cref{subsec:robustness} and \cref{subsec:faq-prompt-engineering}.) To examine if chain-of-thought prompting in this format can reliably trigger effective reasoning across various math word problem datasets, we employed this single set of eight chain-of-thought examples for all benchmarks except AQuA, which is a multiple-choice dataset rather than free response. For AQuA, we utilized four examples with solutions from the training set, as provided in Appendix \cref{tab:appendix-aqua-prompt}.

\textbf{Language models.} We assess five large language models. The first is \textbf{GPT-3} \citep{brown2020language}, where we use the variants text-ada-001, text-babbage-001, text-curie-001, and text-davinci-002, which are believed to correspond to InstructGPT models of sizes 350M, 1.3B, 6.7B, and 175B parameters respectively \citep{ouyang2022training}. The second is \textbf{\lamda{}} \citep{thoppilan2022lamda}, offering models with 422M, 2B, 8B, 68B, and 137B parameters. The third is \textbf{\palm{}}, with models sized 8B, 62B, and 540B parameters. The fourth is \textbf{UL2 20B} \citep{tay2022unifying}, and the fifth is \textbf{Codex} \citep[][code-davinci-002 in the OpenAI API]{chen2021evaluating}. We generate outputs using greedy decoding (although subsequent research suggests that chain-of-thought prompting performance can be enhanced by selecting the majority final answer from multiple sampled outputs \citep{wang2022self}). For \lamda{}, results are averaged over five random seeds, each with a different random permutation of the example order. Since \lamda{} exhibited limited variability across seeds, to conserve computational resources, results for other models are reported using a single exemplar order.

\subsection{Results}
The most significant findings regarding chain-of-thought prompting are illustrated in \cref{fig:main-math}, with comprehensive experimental results for each model family, model scale, and benchmark detailed in \cref{tab:all-lm-math} in the Appendix. There are three main insights. First, \cref{fig:main-math} demonstrates that chain-of-thought prompting is an emergent capability that depends on model size \citep{wei2022emergent}. Specifically, chain-of-thought prompting does not enhance performance for smaller models, and only leads to improved results when applied to models with approximately 100 billion parameters. Our qualitative analysis revealed that smaller models tend to generate fluent but logically flawed chains of reasoning, resulting in worse performance compared to standard prompting.

\input{main_fables/main-math}
Second, chain-of-thought prompting provides greater performance improvements on more complex tasks. For example, on GSM8K (which has the lowest baseline performance), accuracy more than doubled for the largest GPT and \palm{} models. Conversely, for SingleOp, the simplest subset of MAWPS requiring only a single reasoning step, improvements were minimal or even negative (see Appendix \cref{tab:mawps-subsets}).

Third, chain-of-thought prompting applied to GPT-3 175B and \palm{} 540B achieves competitive performance relative to previous state-of-the-art methods, which usually involve fine-tuning a model on a labeled training set. \cref{fig:main-math} illustrates how \palm{} 540B, using chain-of-thought prompting, attains new state-of-the-art results on GSM8K, SVAMP, and MAWPS (though it should be noted that standard prompting had already surpassed the previous best on SVAMP). For the other two datasets, AQuA and ASDiv, \palm{} with chain-of-thought prompting comes within 2\% of the leading performance (see Appendix \cref{tab:all-lm-math}).

To gain a deeper understanding of why chain-of-thought prompting is effective, we conducted a manual inspection of the chains of thought generated by \lamda{} 137B for the GSM8K dataset. Among 50 randomly selected instances where the model produced the correct final answer, nearly all the chains of thought were both logically and mathematically sound, with the exception of two cases where the correct answer was coincidentally reached (refer to \cref{subsec:correct-chain-of-thought} and \cref{tab:appendix-gsm8k-correct-examples} for samples of accurate model-generated chains of thought). Additionally, we reviewed 50 random samples where the model's answer was incorrect. Our analysis revealed that 46\% of these chains of thought were almost correct, containing only minor errors such as calculator mistakes, symbol mapping errors, or a single missing reasoning step, while the remaining 54\% exhibited significant errors related to semantic comprehension or logical coherence (see \cref{subsec:incorrect-chain-of-thought-analysis}).

To shed light on why scaling enhances chain-of-thought reasoning capabilities, we carried out a comparable error analysis for \palm{} 62B and examined how these errors were resolved by scaling up to \palm{} 540B. The findings suggest that increasing \palm{}’s size to 540B corrects a substantial portion of missing one-step errors and semantic understanding mistakes present in the 62B model (see \cref{subsec:why-scale-helps}).

\subsection{Ablation Study}\label{subsec:math-ablation}

The observed advantages of employing chain-of-thought prompting naturally lead to the question of whether similar performance gains can be achieved through alternative prompting strategies. \cref{fig:bar_ablation} presents an ablation study exploring three different variants of chain-of-thought prompting, outlined below.

\textbf{Equation only.} One hypothesis for why chain-of-thought prompting is beneficial is that it generates the mathematical equation to be solved. To test this, we experiment with a variant where the model is instructed to output only the mathematical equation before providing the final answer. As shown in \cref{fig:bar_ablation}, prompting the model to produce only the equation does not significantly improve performance on GSM8K, suggesting that the semantics of GSM8K problems are too complex to be directly converted into an equation without the intermediate natural language reasoning steps found in chain-of-thought prompting. However, for datasets consisting of one-step or two-step problems, we observe that equation-only prompting does enhance performance, as the equation can be more straightforwardly derived from the question (see Appendix \cref{tab:ablations-arithmetic}).

\input{main_fables/ablation-bar}  
\textbf{Compute variation only.} One hypothesis is that chain of thought enables the model to allocate more computational effort (i.e., intermediate tokens) to more challenging problems. To separate the impact of variable computation from the reasoning process in chain-of-thought, we evaluate a setup where the model is prompted to produce a sequence of dots ($\ldots$) matching the number of characters in the equation required to solve the problem. This version performs roughly the same as the baseline, indicating that variable computation alone does not account for the effectiveness of chain-of-thought prompting, and that there seems to be an advantage in explicitly representing intermediate steps through natural language.

\textbf{Chain of thought after the answer.} Another possible advantage of chain-of-thought prompting is that it might help the model more effectively retrieve relevant knowledge learned during pretraining. To test this, we try an alternative approach where the chain-of-thought prompt is provided only after the model has given the answer, thereby determining whether the model truly relies on the generated chain of thought to produce the final answer. This variant performs similarly to the baseline, suggesting that the step-by-step reasoning encapsulated in the chain of thought provides benefits beyond merely triggering knowledge access.

\input{main_fables/main-robustness}  
\subsection{Robustness of Chain of Thought}\label{subsec:robustness}

Sensitivity to the choice of exemplars is a critical factor in prompting methods—for example, changing the order of few-shot examples can cause GPT-3’s accuracy on SST-2 to vary dramatically from near random guessing (54.3\%) to near state-of-the-art performance (93.4\%) \citep{zhao2021calibrate}. In this final section, we examine the robustness of chain-of-thought prompting when using chains of thought authored by different annotators. Besides the results above, which utilized chains of thought written by Annotator A, two other co-authors (Annotators B and C) independently generated chains of thought for the same few-shot examples (shown in \cref{sec:appendix-alternate-annotators}).  
Annotator A also created a more concise version of the chain of thought following the style of solutions presented in \citet{cobbe2021training}.\footnote{For example, while the original chain of thought contains several short sentences (\textit{``There were originally 9 computers. For each of 4 days, 5 more computers were added. So 5 * 4 = 20 computers were added. 9 + 20 is 29.''}), the concise chain reads \textit{``5 * 4 = 20 new computers were added. So there are 9 + 20 = 29 new computers in the server room now.''}}

\cref{fig:bar-robustness-analysis} presents these findings for \lamda{} 137B on the GSM8K and MAWPS datasets (additional ablation results for other datasets are available in Appendix \cref{tab:ablations-arithmetic} / \cref{tab:ablations-commonsense-symbolic}). Although there is variability across different chain of thought annotations, as expected when using exemplar-based prompting \citep{le-scao-rush-2021-many,reynolds2021prompt,zhao2021calibrate}, all sets of chain of thought prompts significantly outperform the standard baseline. This suggests that the effectiveness of chain of thought prompting does not rely on a specific linguistic style.

To verify that successful chain-of-thought prompting generalizes to other exemplar sets, we also conducted experiments using three sets of eight exemplars randomly drawn from the GSM8K training set, which is an independent source (examples in this dataset already contain reasoning steps resembling a chain of thought).\footnote{We select examples $\leq 60$ tokens to fit within our input context window, and restrict to examples with $\leq 2$ steps to enable a fair comparison with the eight exemplars we composed.} As shown in \cref{fig:bar-robustness-analysis}, these prompts performed similarly to our manually crafted exemplars, and also notably outperformed standard prompting.

Beyond demonstrating robustness to different annotators, independently created chains of thought, varying exemplars, and diverse language models, we also observe that chain-of-thought prompting for arithmetic reasoning maintains effectiveness across different exemplar orders and varying numbers of exemplars (refer to \cref{subsec:faq-prompt-engineering}).

\section{Commonsense Reasoning}\label{sec:commonsense-reasoning}

While chain of thought prompting is especially well-suited for math word problems, its language-based framework makes it applicable to a wide range of commonsense reasoning tasks, which require reasoning about physical and human interactions assuming general background knowledge. Commonsense reasoning is essential for engaging with the world and remains beyond the capabilities of current natural language understanding systems \citep{talmor2022commonsenseqa}.

\textbf{Benchmarks.}  
We utilize five datasets that encompass a wide variety of commonsense reasoning challenges. The widely-used \textbf{CSQA} \citep{talmor-etal-2019-commonsenseqa} involves answering commonsense questions about the world, requiring intricate semantics and often background knowledge. \textbf{StrategyQA} \citep{geva-etal-2021-aristotle} demands that models deduce a multi-step strategy to respond to questions. From the BIG-bench initiative \citep{bigbench}, we select two specialized evaluation sets: \textbf{Date} Understanding, which entails inferring a date from a given context, and \textbf{Sports} Understanding, which involves judging whether a sports-related sentence is plausible or not. Lastly, the \textbf{SayCan} dataset \citep{ahn2022can} requires converting natural language instructions into a sequence of robot actions chosen from a discrete set.  
\cref{fig:dataset-examples} provides sample examples with chain of thought annotations for each dataset.

\textbf{Prompts.}  
Our experimental setup aligns with the prior section. For CSQA and StrategyQA, we randomly picked samples from the training sets and manually crafted chains of thought to serve as few-shot exemplars. Since the two BIG-bench tasks lack training sets, we selected the initial ten examples from their evaluation sets as few-shot exemplars and report results on the remaining evaluation examples. For SayCan, we employ six examples from the training set used in \citet{ahn2022can}, for which we also manually created chains of thought.

\textbf{Results.}
\cref{fig:commonsense-results} presents these findings for \palm{} (complete results for \lamda{}, GPT-3, and various model sizes can be found in \cref{tab:all-lm-commonsense}).
Across all tasks, increasing the model size enhanced performance with standard prompting; applying chain-of-thought prompting yielded additional improvements, with the largest gains observed for \palm{} 540B. Using chain-of-thought prompting, \palm{} 540B delivered strong results compared to baselines, surpassing the previous state of the art on StrategyQA (75.6\% vs 69.4\%) and outperforming an unaided sports enthusiast on sports knowledge (95.4\% vs 84\%).
These outcomes indicate that chain-of-thought prompting can also boost performance on tasks involving diverse commonsense reasoning skills (although the improvement was minimal for CSQA).

\input{main_fables/main-commonsense}

\input{main_fables/main-symbolic}
\section{Symbolic Reasoning}\label{sec:symbolic-reasoning}
Our final set of experiments focuses on symbolic reasoning, which is straightforward for humans but potentially difficult for language models. We demonstrate that chain-of-thought prompting not only enables language models to handle symbolic reasoning problems that are challenging under standard prompting but also supports length generalization to test inputs that are longer than the few-shot examples provided.

\paragraph{Tasks.}
We consider the following two simple tasks.

\begin{itemize}[leftmargin=*,topsep=0pt]
    \itemsep0em \item \textbf{Last letter concatenation.} This task requires the model to concatenate the last letters of each word in a name (e.g., \textit{``Amy Brown''} $\rightarrow$ \textit{``yn''}). 
    It represents a more difficult variant of first letter concatenation, which language models can already solve without chain-of-thought prompting.\footnote{We tested 10 common names using GPT-3 \texttt{davinci} and it answered all but one correctly.} Full names are generated by randomly combining names from the top one-thousand first and last names based on name census data ({\small{\url{https://namecensus.com/}}}).
    \item \textbf{Coin flip.} In this task, the model must determine whether a coin remains heads up after a sequence of people either flip or do not flip it (e.g., \textit{``A coin is heads up. Phoebe flips the coin. Osvaldo does not flip the coin. Is the coin still heads up?''} $\rightarrow$ \textit{``no''}).
\end{itemize}

Since the construction of these symbolic reasoning tasks is clearly defined, for each task we consider an \textit{in-domain} test set where examples have the same number of steps as the training or few-shot exemplars, as well as an \textit{out-of-domain} (OOD) test set where evaluation examples involve more steps than those in the exemplars. For the last letter concatenation task, the model is only exposed to exemplars of names with two words, and then it performs last letter concatenation on names consisting of 3 and 4 words.\footnote{For names longer than 2 words, we concatenate multiple first and last names together.} Similarly, we vary the number of potential flips in the coin flip task. Our experimental setup employs the same methods and models as described in the previous two sections. We once again manually construct chains of thought for the few-shot exemplars for each task, as shown in \cref{fig:dataset-examples}.

\paragraph{Results.}
The outcomes of these in-domain and OOD assessments are presented in \cref{fig:symbolic-main} for \palm{}, with results for \lamda{} provided in Appendix \cref{tab:all-lm-symbolic}. Using \palm{} 540B, chain-of-thought prompting achieves nearly 100\% success rates (note that standard prompting already solves the coin flip task with \palm{} 540B, although this is not the case for \lamda{} 137B). It is important to observe that these in-domain tests are considered "toy tasks" since the exact solution steps are already supplied by the chains of thought in the few-shot exemplars; the model simply needs to replicate these steps with new symbols in the test examples. Despite this, smaller models still fail— the capability to perform abstract manipulations on unseen symbols in these three tasks emerges only at the scale of models with around 100B parameters.

Regarding the OOD evaluations, standard prompting fails on both tasks. When using chain-of-thought prompting, language models display upward scaling trends (although performance remains lower compared to the in-domain scenario). Thus, chain-of-thought prompting enables length generalization beyond previously encountered chains of thought for sufficiently large language models.

\section{Discussion}\label{sec:discussion}

We investigated chain-of-thought prompting as a straightforward approach to induce multi-step reasoning in large language models. Initially, we observed that chain-of-thought prompting significantly boosts performance on arithmetic reasoning, delivering gains far exceeding those of ablations and demonstrating robustness across various annotators, exemplars, and language models (\cref{sec:arithmetic-reasoning}). Subsequently, experiments on commonsense reasoning highlighted that the linguistic structure of chain-of-thought reasoning makes it broadly applicable (\cref{sec:commonsense-reasoning}). Finally, we demonstrated that in symbolic reasoning, chain-of-thought prompting supports OOD generalization to longer sequences (\cref{sec:symbolic-reasoning}). Throughout all experiments, chain-of-thought reasoning was induced solely by prompting off-the-shelf language models, with no finetuning performed during this study.

The emergence of chain-of-thought reasoning correlated with model scale has been a consistent finding \citep{wei2022emergent}. For many reasoning tasks where standard prompting yields a flat scaling curve, chain-of-thought prompting results in sharply rising scaling curves. Chain-of-thought prompting seems to broaden the spectrum of tasks that large language models can handle effectively—indicating that standard prompting only sets a lower bound on large language models’ capabilities. This insight likely prompts more questions than answers—for example, how much further might reasoning skills improve with increased model scale? What other prompting techniques could extend the range of solvable tasks for language models?

Regarding limitations, we first clarify that although chain of thought mirrors human reasoning processes, this does not confirm that the neural network is genuinely “reasoning,” which remains an open question. Second, while the cost of manually adding chain-of-thought exemplars is minimal in a few-shot setting, such annotation expenses might be prohibitive for finetuning (though this might be addressed by synthetic data creation or zero-shot generalization). \orange{Third, there is no assurance that the reasoning paths are correct, which can yield both accurate and inaccurate answers; improving the factual accuracy of language model outputs remains a direction for future research \citep[][\textit{inter alia}]{rashkin2021measuring,ye2022unreliability,wiegreffe2021reframing}.} Finally, since chain-of-thought reasoning emerges only in large-scale models, it is expensive to deploy in practical applications; further investigation could focus on methods to induce reasoning in smaller models.

\section{Related Work} This study draws inspiration from various research fields, which we elaborate on in an extended related work section (\cref{sec:extended-related-work}). Here, we highlight two key directions and their associated studies that are particularly pertinent.

The first key area involves employing intermediate steps to address reasoning challenges. \cite{ling-etal-2017-program} introduced the concept of using natural language rationales to tackle math word problems through a sequence of intermediate stages. Their approach contrasts notably with research that utilizes formal languages for reasoning \citep{roy-2016-reasoning, chiang-chen-2019-semantically, amini-etal-2019-mathqa, chen2019neural}. \citet{cobbe2021training} build upon \citet{ling-etal-2017-program} by developing a larger dataset and using it to fine-tune a pretrained language model instead of training a model from the ground up. In program synthesis, \cite{nye2021show} employ language models to predict the final results of Python programs by first predicting intermediate computational outcomes line-by-line, demonstrating that this stepwise prediction approach outperforms directly predicting the final outputs.

This paper is also closely connected to the extensive recent research on prompting. Since the introduction of few-shot prompting by \citet{brown2020language}, various general strategies have enhanced the prompting capabilities of models, such as automatically learning prompts \citep{lester-etal-2021-power} or providing models with instructions that describe a task \citep{wei2021finetuned,sanh2021multitask,ouyang2022training}. While these methods focus on improving or supplementing the input portion of the prompt (e.g., instructions prepended to inputs), our approach takes a different route by enriching the outputs of language models with a chain of thought.

\section{Conclusions} We have investigated chain-of-thought prompting as a straightforward and widely applicable technique for boosting reasoning capabilities in language models. Through experiments involving arithmetic, symbolic, and commonsense reasoning, we observe that chain-of-thought reasoning emerges as a property related to model scale, enabling sufficiently large language models to tackle reasoning tasks that otherwise exhibit flat scaling performance. Expanding the variety of reasoning tasks that language models can handle may encourage further research into language-based reasoning methods.

\end{document}